% Native mathematical source; requires amsmath. Integrate tags through the author template.
\begin{align*}
  i & = x-A_x,         & j & = s_y(y-A_y), \\
  e & = d_x-D-Di+d_xj.
\end{align*}
An $x$ step subtracts $D$ from $e$; a $y$ step adds $d_x$.
For an $x$-major line with $d_x>0$, the next minor coordinate steps forward when
\begin{align*}
  \frac{D(i+1)}{d_x}-j & \ge \frac12
                       & \Longleftrightarrow\quad 2e & \le d_x.
\end{align*}
For a $y$-major line with $D>0$, the corresponding condition is
\begin{align*}
  \frac{d_x(j+1)}{D}-i & \ge \frac12
                       & \Longleftrightarrow\quad 2e & \ge -D.
\end{align*}
Both decisions use the same pre-update value of $2e$. Equality chooses the forward
minor step after endpoint normalization. The coincident-point case exits before
these positive-denominator arguments are needed.


\paragraph{Major-axis progress.}
On an $x$-major line, let $\delta=j-Di/d_x$. Nearest rounding gives
$|\delta|\le\tfrac12$, hence
\[
  e=d_x-D+d_x\delta\ge\frac{d_x}{2}-D\ge-\frac D2
  \qquad(d_x\ge D).
\]
Thus the $x$ decision always succeeds before the endpoint.
On a $y$-major line, let $\delta=i-d_xj/D$. Then
\[
  e=d_x-D-D\delta\le d_x-\frac D2\le\frac{d_x}{2}
  \qquad(D\ge d_x),
\]
so the $y$ decision always succeeds. The minor-axis test gives the next rounded
sample. Starting at $i=j=0$, induction preserves that relation and reaches the
endpoint after $\max(d_x,D)$ updates. Painting before the endpoint test gives
$\max(d_x,D)+1$ samples, at most eight in this assignment. A point exits without
an update; each division in this argument is used only for a nonzero major extent.
